\begin{corollary}[Proved under stated hypotheses: the three-point surface in every genus]
\label{D:schroer-all-genera}
In Schr\"oer's construction of a normal proper surface $S$ with
$\operatorname{Pic}(S)=0$, let the initial smooth proper curve
$C$ have any positive genus, and let $s_0,s_1,s_2$ be the three
contracted points.  Then
\[
                  \pi_1(S\setminus\{s_0,s_1,s_2\})
\]
is topologically finitely generated.  In particular, the
higher-genus version of this explicit nonprojective surface is
not a counterexample to D.2(ii).
\end{corollary}
\begin{proof}
Characteristic zero is already covered by
Proposition~\ref{D:fixed-boundary-positive}; assume that the
characteristic is positive.  Use the notation and point choices
of \cite[\S\S2.5 and 3.1--3.2, author's PDF pp.~4--5]{Schroer1999D}.
Let $X$ be the blowup of $\mathbb P^1\times C$ at
$(b_0,c_1),(b_1,c_1),(b_2,c_2)$, let $R_i$ be the strict
vertical fibres, and let $A$ be the strict horizontal curve
through $c_1$.  The common open is
\[
 U=X\setminus(R_0\cup R_1\cup A\cup R_2)
   =S\setminus\{s_0,s_1,s_2\}.
\]

We identify the formal neighborhood of $R_i$ for $i=0,1$.
Choose a coordinate $t$ on an affine neighborhood of $b_i$
in $\mathbb P^1$, with $t(b_i)=0$, and a parameter $z$ at
$c_1$ on $C$.  The blowup chart containing the strict transform
of $\{b_i\}\times C$ has coordinates
\[
                t=zu,\qquad R_i=\{u=0\}.
\]
Away from $c_1$ the normal coordinate is $t$.  On the overlap
the transition is $u=t/z$.  These are exactly the linear
transition functions of the line bundle
$\mathcal O_C(-c_1)$.  They identify the completed neighborhood
of $R_i$ with the completion of its zero section.  The other
two blowup centres are disjoint from $R_i$ and do not change
this formal neighborhood.  The line bundle has degree $-1$.
Proposition~\ref{D:graded-cone-tame} therefore shows that any
finite cover of the punctured germ at $s_i$ is tame at the
divisorial valuation $R_i$, including separability of its
residue extension.

Let $V\to U$ be a connected finite \'etale cover.  Its finite
normalization over $S$ is \'etale away from the three omitted
points.  Applying the preceding local calculation at $s_0$
and $s_1$ proves tameness at $R_0$ and $R_1$ for this cover.
There is no restriction on its degree, and no restriction on
the genus of $C$ in this calculation.

Section 3.1 of the cited source constructs a normal projective
surface $T$ by contracting only $A\cup R_2$.  The construction
works for the same arbitrary positive-genus curve; it is the
partial contraction used in Corollary~\ref{D:schroer-elliptic}.
If $q$ is the image of $A\cup R_2$, then
\[
                   U=T\setminus(R_0\cup R_1\cup\{q\}).
\]
The only divisorial boundary components on $T$ are $R_0$ and
$R_1$.  We have proved that every finite \'etale cover of $U$
is tame at both.  The omitted point $q$ has codimension two,
so it adds no condition to this fixed-boundary tame category.
Consequently
\[
                   \pi_1(U)=Q(T,T\setminus U).
\]
The right-hand side is topologically finitely generated by
Proposition~\ref{D:fixed-boundary-positive}, because $T$ is
projective.  We have imposed no log-canonical or formal-cone
condition on the third singularity $q$.
\end{proof}
